\section{Introduction}
Two pelleting operations can consume different amounts of electricity for entirely ordinary reasons: one processes more material, uses a different product formulation or runs on another machine. These differences matter when a plant tries to judge whether an operation consumed more energy than expected. A single kWh/t reference is convenient, but it can leave much of the variation in production conditions unexplained. An energy baseline should account for that variation while remaining clear enough for engineers to understand why its estimate changes.

The physical basis for such adjustment is well established. Feed properties and processing conditions jointly influence pellet quality and manufacturing behaviour \citep{thomas1996,thomas1997}. Conditioning changes the thermal and moisture state of the material before compaction \citep{thomas2020}, and both steam-conditioning rate and feed composition can affect the process and its mechanical energy demand \citep{skoch1981,bastiaansen2023}. Industrial measurement and verification must nevertheless distinguish changes in production from savings attributable to an intervention \citep{kissock2008}. Related work on automated energy baselines shows why predictive accuracy and uncertainty matter for that distinction \citep{granderson2016,gallagher2018}.

Whether operating histories improve on production-adjusted estimates remains an empirical question. Regularized regression and boosting can capture associations between production inputs and consumption \citep{hoerl1970,friedman2001}; temporal architectures such as TFT and N-HiTS can also use operating histories \citep{lim2021,challu2023}. Pretrained models, including TimesFM and Chronos, offer another route when local training data are limited \citep{das2024,ansari2024}. Their usefulness depends on what the industrial records actually represent. A sequence of completed operations, each with a different mass and duration, is not equivalent to a regularly sampled power signal. The prediction origin and the information available at that time must therefore be defined before comparing models \citep{hewamalage2023}.

Uncertainty is equally consequential. An estimate can be close on average and still miss particular machines or periods. Prediction intervals help express this variability, provided that coverage is considered together with interval width \citep{gneitingcal2007,gneiting2007}. Conformal calibration provides a practical starting point \citep{angelopoulos2021}, although temporal dependence and changing operating conditions complicate its guarantees \citep{barber2023}. Good coverage over a complete dataset also need not imply adequate coverage within each machine or product group \citep{barber2021}.

Here we use real records from three feed-pelleting lines to develop an explainable baseline for electrical energy per operation. We first compare production-adjusted models with a machine-specific energy-intensity reference and assess the selected model in a chronological test. We then ask whether completed same-machine histories improve estimation through linear, neural or pretrained temporal models. Finally, we examine the baseline's variable contributions and the scope for updating its point estimates and uncertainty using recent residuals. The aim is a reproducible statistical component that can later be incorporated into a digital twin, with a clear account of the evidence needed for prospective use. Figure~\ref{fig:framework} illustrates the selected baseline and shows how its output is read for an actual industrial record.

\fig{00_framework}{framework}{1}{The selected energy baseline. (a) Recorded production conditions are transformed into an electricity estimate; residuals from a later calibration period determine its uncertainty limits. (b) A real operation on Pellet~2 illustrates the output. The light segment is the fixed 95\% interval and the dark segment the 90\% interval; the circle marks the estimate and the diamond the observed energy. The example retains source row 2080 and uses its archived values. An excess above the point estimate is a model residual, not a measured saving or an adjudicated fault.}
